\documentclass{article}

% Language setting
% Replace `english' with e.g. `spanish' to change the document language
\usepackage[english]{babel}

% Set page size and margins
% Replace `letterpaper' with `a4paper' for UK/EU standard size
\usepackage[letterpaper,top=2cm,bottom=2cm,left=3cm,right=3cm,marginparwidth=1.75cm]{geometry}

% Useful packages
\usepackage{amsmath}
\usepackage{graphicx}
\usepackage[colorlinks=true, allcolors=blue]{hyperref}

\title{Al Medina}
\author{Website onderzoek}

\begin{document}
\maketitle

\pagebreak

\tableofcontents

\pagebreak

\section*{Inleiding}
    
Dit vooronderzoeksdocument is opgesteld ter voorbereiding op de ontwikkeling van de nieuwe website voor Al Medina. Het doel van dit document is om een helder overzicht te geven van de vereisten, verwachtingen en doelen die met de website bereikt moeten worden. Door middel van gedetailleerd onderzoek naar het bedrijf, de doelgroep, en de concurrentie, biedt dit document een solide basis voor de creatie van een effectieve en gebruiksvriendelijke website die aansluit bij de bedrijfsstrategie.

De nieuwe website zal niet alleen dienen als visitekaartje voor het bedrijf, maar ook als een krachtig hulpmiddel om de online aanwezigheid te versterken, de interactie met klanten te verbeteren, en de bedrijfsdoelen te ondersteunen. Dit document behandelt alle aspecten die nodig zijn voor de succesvolle ontwikkeling van de website, waaronder technische vereisten en designvoorkeuren.

Het doel van dit vooronderzoek is om ervoor te zorgen dat de website niet alleen voldoet aan de huidige behoeften van Al Medina, maar ook flexibel genoeg is om in de toekomst mee te groeien met het bedrijf en de veranderende markt.

\pagebreak

\section{Bedrijfsinformatie}
\subsection{Naam van het bedrijf}
Al Medina Cupping
\subsection{Bedrijfsdoelen}
\begin{enumerate}
    \item Informatie over diensten geven
    \item Meer klanten bereiken
    \item Meer afspraken krijgen
    \item Optimaliseren van bedrijfsprocessen
\end{enumerate}

\subsection{Branche}
Het bedrijf opereert binnen de alternatieve geneeskunde sector.


\section{Doelgroepanalyse}
\subsection{Demografie}
De doelgroep van Al Medina is volwassen vrouwen en kinderen in de omgeving van Purmerend van alle afkomsten.

\subsection{Behoeften en doelstellingen}
Potentiële bezoekers van de website willen:
\begin{enumerate}
    \item Een afspraak boeken om via alternatieve geneesmiddelen hun lichamelijke klachten te verminderen;
    \item Zich inlezen in de diensten van Al Medina om er achter te komen of de deze voor hen bestemd zijn.
\end{enumerate}

\section{Concurrentieanalyse}
\subsection{\href{https://www.skinandmore.nl/}{Skin and more}} 

Sterke punten:
\begin{enumerate}
    \item Mooi ontwerp
    \item Pakkende en nette headers
    \item Goede animaties
\end{enumerate}

\noindent Zwakke punten:
\begin{enumerate}
    \item Navigatiebalk beweegt te veel
    \item Sommige animaties gaan fout
    \item Slechte performance score (61)
\end{enumerate}

\subsection{\href{https://bodytemplecupping.nl/}{Body Temple Cupping}}

Sterke punten:
\begin{enumerate}
    \item Goede inrichting homepagina
    \item Duidelijke CTA's op homepagina
    \item F.A.Q. sectie is erg behulpzaam
    \item Professionele uitstraling
\end{enumerate}

\noindent Zwakke punten:
\begin{enumerate}
    \item Sommige functies (instagram feed) werken niet.
    \item Contact in de footer bevat weinig CTA's.
\end{enumerate}

\subsection{\href{https://hijamaencupping.nl/}{Hijama en Cupping}}

Sterke punten:
\begin{enumerate}
    \item Goed gebruik van kleurenschema
    \item Moderne uitstraling met mooie kopjes
    \item Goede onderscheiding tussen secties
\end{enumerate}

\section{Functionaliteiten}


\end{document}